
\documentclass[final]{elsarticle}

\usepackage[margin=2.50cm]{geometry}  % tweak margins
%\usepackage{draftcopy}                % uncomment to print a "DRAFT-COPY" in background
\usepackage{graphicx}                 % if you want to include PostScript figures
\usepackage[figuresright]{rotating}   % if you have landscape tables
\usepackage{multirow}                 % fancy tables
\usepackage[toc,page]{appendix}       % appendix
\usepackage{lineno}                   % for printing the line numbers
\usepackage{natbib}                   % for building bibliography with \cite{}
\usepackage{setspace}
\usepackage{times}
\usepackage{afterpage}
\usepackage{float}
\usepackage{subcaption}
\usepackage{longtable}
\usepackage{blkarray}                 % for correlation matrices
\usepackage{color}
\usepackage{rotate}                   % rotate text in tables
\usepackage{lscape}                   % allow page to be landscape for large tables
\usepackage{amsfonts}                 % extra symbols
\usepackage{amsmath}
%\usepackage{floatpag}                 % \thisfloatpagestyle{empty}, to suppress page numbers
\usepackage{lmodern}
\usepackage{dcolumn}% Align table columns on decimal point
\usepackage{caption}% Use captionof to be able to minipage figure's \& tables together, but still get correct captions
\usepackage{enumitem}

\usepackage[breaklinks=true,hyperfigures=true]{hyperref} % make contents and references links, but do not colour them

\setenumerate{listparindent=\parindent} % Use paragraph indentation within the enumerate

\begin{document}

\begin{enumerate}

\item 

\textcolor{blue}
{To all three TN authors:
could you please discuss consistency of your analysis results with others?
For example, the p-theta and Valor can compare their p-values, especially 
by using rate-only p-values.  
According to the Kirsty's presentation at 18June15 analysis meeting, Valor and MaCh3
can be compared by using likelihood ratio Bayes Factor.  Such information should also be
included in the technotes as a record.
It is also nice if p-theta can be compared with MaCh3 as well (although I don't know it is possible or not).}

DONE. Added a short section at the end of the technical note.

\end{enumerate}

\end{document}














